

Multimodal AI systems are moving into real-world use. They help diagnose diseases, moderate content, and assist drivers. This growing deployment raises a pointed question: \emph{can we trust what they see?} Imagine a radiologist reviewing a chest X-ray with an AI assistant. The system receives both the image and the radiologist's preliminary notes. If those notes contain an error, will the AI catch it by looking at the image? Or will it simply echo the mistake?

Unfortunately, the evidence points to the second outcome. Show a model an image of a red car alongside text claiming ``the car is blue,'' and most MLLMs will confidently answer ``blue''~\cite{Insight_Over_Sight,PhD_ChatGPT,wu2024autohallusion}. They
trust the text over their own visual perception. This pattern appears
across different model families~\cite{llava-1.5,instructblip2023,bai2025qwen2,bai2025qwen3vltechnicalreport,chen2024internvl} and benchmarks~\cite{zhang2025robust,jia2025benchmarking,qian2024easy}. In safety-critical applications~\cite{bannur2024maira,hou2025one}, such blind trust in text could cause serious harm.

\begin{figure}[t]
  \centering
  \includegraphics[width=\linewidth]{figures/teaser.pdf}
  \caption{How predictions shift under visual-textual conflict. (a) The model sees a red car but outputs the text-suggested answer. (b) Modal Dominance Ratio (MDR) across layers: early layers favors vision, but late layers override toward text. (c) Shift direction predicts correctness—text-ward shifts correlate with failures.}
  \label{fig:teaser}
\end{figure}

\emph{Why} does this happen? Prior work describes the \textit{\underline{behavior}} but not the \textit{\underline{mechanism}}~\cite{Insight_Over_Sight,zhang2025robust}. Is visual information poorly encoded in the first place? Is it lost during cross-modal fusion? Or is it correctly processed but later overwritten? Understanding the failure mode matters: different causes require different solutions.


To find out, we trace how a model's predictions evolve across its layers. Following~\cite{depth_adaptive2020,confident_adaptive2022}, we project hidden states at each layer to output probabilities. This lets us observe what the model ``prefers" at each depth, and what we find is surprising. Surprisingly, in many conflict failures, \textbf{the model assigns higher probability to the correct visual answer in intermediate layers}. It then reverses course and outputs the text answer (Figure~\ref{fig:teaser}(a-b)).


We term this phenomenon \textbf{late-layer textual override}. To measure it, we introduce
Modal Dominance Ratio (MDR), which quantifies vision-versus-text
preference at each layer. In failure cases, MDR starts positive (visual)
and flips negative (textual) in later layers. The visual information was there; it simply didn't survive the full forward pass.

\textbf{Not every late-layer shift is harmful.}
Such prediction shift occurs in successful cases too~\cite{wang2025shift}, not just failures. What distinguishes the two outcomes is the \emph{direction} of the shift. We quantify this using the change in MDR across the transition point. As shown in Figure~\ref{fig:teaser}(c), in failures, 85\% of transitions move toward text. In successes, 89\% move toward vision. Late layers sometimes help; they sometimes hurt. The key
is knowing which.

This directional asymmetry points toward a practical solution. If a model holds a confident visual prediction at the transition layer but abandons it by the final layer, we have reason to suspect harmful override. We capture this intuition with two signals derived from our analysis. \textit{Anchor confidence} measures how certain the transition-layer prediction is. \textit{Prediction retention} captures whether it survives to the output. A high anchor confidence paired with low retention is the signature of harmful override.

\textbf{Recovering what was known.} These observations motivate \textbf{\underline{C}}onflict-\textbf{\underline{A}}ware \textbf{\underline{L}}ayer \textbf{\underline{R}}eference \textbf{\underline{D}}ecoding (CALRD), a training-free method that restores intermediate predictions when override is detected. Our core insight is that predictions at the transition layer still carry the vision-grounded answer before it gets overwritten. By blending transition-layer logits back into the final distribution, we can recover it. CALRD first locates the layer where the output distribution shifts most sharply. It then uses anchor confidence and prediction retention to set correction strength: when both signals point to harmful override, CALRD blends in the transition-layer logits; otherwise, it leaves the output unchanged.


CALRD is straightforward. The model already has the right answer in its intermediate layers; we are helping it remember. No external knowledge or modules are involved. This points to a shift in perspective: improving MLLM reliability may be less about teaching models to see better and more about helping them retain what they have already seen.


We evaluate CALRD on five MLLMs spanning architectures, from InstructBLIP~\cite{instructblip2023} to Qwen3-VL~\cite{bai2025qwen3vltechnicalreport}. Results show consistent improvements on conflict benchmarks  (Conflict-VQA, PhD-icc~\cite{PhD_ChatGPT}) while maintaining performance on standard hallucination evaluations (POPE~\cite{pope2023}, CHAIR\cite{Object_Hallucination}). The effectiveness across model families suggests that late-layer override is a general phenomenon, not an artifact of specific architectures. Our contributions are as follows:


\begin{itemize}[leftmargin=*,nosep]
 \item We introduce Modal Dominance Ratio to trace layer-wise  modal preference and uncover late-layer textual override as a characteristic failure pattern. We show that transition direction,  not mere occurrence, predicts correctness.
 \item We propose CALRD, a training-free method that detects harmful overrides through complementary signals and applies adaptive correction, preserving beneficial processing.
 \item Experiments show that CALRD achieves 4--9\% absolute gains on conflict tasks without degrading standard performance. We also contribute Conflict-VQA, a diagnostic benchmark with explicit visual-textual answer annotations for mechanistic analysis.
\end{itemize}


\begin{figure}[t]
\centering
  \includegraphics[width=0.99\linewidth]{figures/data.pdf}
\caption{Conflict-VQA construction pipeline. We process images from VrR-VG and questions from TDIUC, use GPT-4 to generate $C_{\text{factual}}$ and $C_{\text{conflict}}$, and verify all samples manually.}
  \label{fig:data_pipeline}
\end{figure}

\begin{table}[t]
\centering
\resizebox{\columnwidth}{!}{
\begin{tabular}{ccccc}
\toprule
InstructBLIP & LLaVA-1.5 & LLaVA-1.6 & Qwen2.5-VL & Qwen3-VL \\
\midrule
2,017 & 1,212 & 2,506 & 3,940 & 5,252 \\
\bottomrule
\end{tabular}
}
\caption{Competent subset size per model. Total samples: 5,969. A sample is competent if the model answers correctly in all five trials under non-conflict context.}
\label{tab:data_stats}
\end{table}
